\documentclass{article}
\usepackage{amsmath, amssymb, mathtools, amsthm, mathrsfs}
\usepackage{tikz-cd}

%\geometry{margin = 1in}
\usepackage[margin=60pt, tmargin=90pt, bmargin=90pt]{geometry}

\newcommand{\N}{\mathbb N}
\newcommand{\Z}{\mathbb Z}
\newcommand{\Q}{\mathbb Q}
\newcommand{\R}{\mathbb R}
\newcommand{\C}{\mathbb C}
\def\->{\ensuremath\rightarrow}
\def\=>{\ensuremath\Rightarrow}
\def\<-{\ensuremath\leftarrow}
\def\<={\ensuremath\Leftarrow}
\newcommand{\set}[1]{\{#1\}}
\newcommand{\gen}[1]{\langle#1\rangle}
\newcommand{\lspace}{\vspace{6pt}}

\title{Hartshorne Chapter 1}
\author{H.C.}
\date{August 2025}

\begin{document}
\maketitle

\section*{Section 1.2}
\subsection*{Question 2.1}

Consider $Z(\mathfrak a)$ instead as a cone in $n+1$-affine space. Then as $\deg f > 0$ by assumption, $Z(f)$ cannot include the origin as it is homogeneous. Hence the behaviour of $Z(\mathfrak a)$ at the origin can be ignored, and it is well-defined everywhere else. But then $f \in \sqrt{\mathfrak a}$ by the affine Nullstellensatz, and we are done. \qed
\lspace

\subsection*{Question 2.2}

(i) $\Leftrightarrow$ (ii):

Consider $Z(\mathfrak a)$ as a cone in affine $n+1$-space. Then $Z(\mathfrak a) = \emptyset$ in $\mathbf{P}^n$ iff $Z(\mathfrak a) = \set{0}$ or $\emptyset$ in $\mathbf{A}^{n+1}$. By the affine Nullstellensatz, $\sqrt{\mathfrak a}$ is going to be $\gen{x_1, \dots, x_{n+1}}$ or $k[x_1,\dots, x_{n+1}]$ respectively, which correspond to the homogeneous ideals $S$ and $S_+ = \bigoplus_{d > 0} S_d$ respectively.

\noindent (ii) $\Leftrightarrow$ (iii):

Forward direction is clear.

Now suppose that $S_d \subseteq \mathfrak a$ for some $d>0$. Then $S_1 \subseteq \sqrt{S_d} \subseteq \sqrt{\mathfrak a}$. Hence $\mathfrak a$ is a homogeneous ideal that contains every positive degree homogeneous polynomial. Depending on if it contains a constant or not, $\sqrt{\mathfrak a} = S$ or $S_+$. \qed
\lspace

\subsection*{Question 2.3}

Skipped.

\subsection*{Question 2.4}

Skipped.

\subsection*{Question 2.5}

Skipped.

\newpage
\subsection*{Question 2.6}

WLOG suppose that we are working with $x_0$ and $U_0$.

Let $\beta: k[x_1,\dots, x_n] \to S$ denote the homogenisation homomorphism as given in Proposition 2.2. Now consider some homogenised polynomial $\beta(f)$. Send $\beta(f)$ to $\frac{\beta(f)}{x_0}$, and let $\phi: k[x_1,\dots,x_n] \to S_{(x_0)}$ denote the composite map. It defines a ring homomorphism that sends $k[x_1,\dots, x_n]$ to the subring of $S_{(x_0)}$ of degree $0$ elements.

Equally consider the evaluation ring homomorphism $\alpha: S \to k[x_1,\dots, x_n]$ given Proposition 2.2., and extend its domain to $S_{(x_0)}$.

To prove that $\phi$ and $\alpha$ define isomorphisms, it suffices to check monomials. Hence let $\frac{g}{x_0^k} \in S_{(x_0)}$ be such a monomial, such that $x_0 \nmid g$. Then $\phi(\alpha(g)) = \phi(g) = \frac{g}{x_0^{\deg g}}$ as desired. Equally if $f \in k[x_1,\dots,x_n]$ is a monomial we have that $\alpha(\phi(f)) = \alpha(\frac{f}{x_0^{\deg f}}) = f$ as desired.

Now note that if $\mathfrak a \lhd S$ is the radical ideal given by $I(Y)$, $A(Y_0)$ is exactly $k[x_1,\dots,x_n] / \alpha (\mathfrak a)$. Hence by considering the isomorphism given above we get that $A(Y_p)$ is isomorphic to the subring of degree $0$ elements in $S(Y)_{(x_0)}$. But then every fraction in $S_{(x_0)}$ can be multiplied by $x_0$ to the appropriate power, positive or negative, such that the result has degree $0$. Hence $S(Y)_{(x_i)} \simeq A(Y_i)[x_i,x_i^{-1}]$ for all $i$.

\noindent Hence, by that result $Q(S(Y)) = Q(S(Y)_{(x_i)}) \simeq Q(A(Y_i)[x_i,x_i^{-1}]) = Q(A(Y_i))(x_i) \simeq Q(S(Y)_{(x_i)}\vert_{\deg = 0})(x_i)$. But then by Theorem 1.8Aa, $\dim(S(Y)) = \text{trdeg}_k (Q(S(Y))) = \text{trdeg}(Q(S(Y)_{(x_i)}\vert_{0})(x_i)) = \text{trdeg}(Q(S(Y)_{(x_i)}\vert_{0})) +1 = \dim(A(Y_i)) + 1 = \dim(Y_i) + 1$.

Moreover, as the $U_i$'s form a finite open cover for $Y$, we have that $\dim(S(Y)) = \dim(Y)+1$. Finally notice that as the equality $\dim(Y)+1 = \dim(S(Y)) = \dim(Y_i) + 1$ must hold for all $i$, as long as $Y_i$ is nonempty such that $\dim(Y_i)$ is well-defined we get that $\dim(Y_i) = \dim(Y)$ necessarily. \qed
\lspace
\lspace

\textit{Note.} Sections of the proof are shamelessly copied from https://williamtroiani.github.io/pdfs/HartshorneSolutions.pdf because there was no way I was figuring out the dimensionality argument.
\lspace

\subsection*{Question 2.7}

Skipped.

\subsection*{Question 2.8}

Notice that a projective variety $Y \subseteq \mathbf{P}^n$ has dimension $n-1$ iff $\dim(S(Y)) = n$ by the result of 2.6 above. But then $\dim(S(Y)) = n$ iff the prime ideal $\mathfrak p:= I(Y)$ necessarily has height one in $k[x_1,\dots,x_n]$ by Theorem 1.8A. But then by Prop 1.12A this is equivalent to $\mathfrak p = \gen{f}$ for some irreducible $f$. Hence $Y$ is a hypersurface iff it has dimension $n-1$. \qed
\lspace

\subsection*{Question 2.9}

\noindent a):

That $\beta(I(Y)) \subseteq I(\overline Y)$ is clear.

Now suppose that $f \in I(\overline Y)$ is a homogeneous polynomial that vanishes on $Y$. Then $\alpha(f) \in I(Y)$, but $\beta(\alpha(f)) = f$ ($\beta \circ \alpha$ is not the identity in general but it is on homogeneous elements). Hence we get equality on the homogeneous terms, such that $I(\overline Y) = \gen{\beta(I(Y))}$ as it is necessarily generated by homogeneous elements. \qed

\noindent b):

From Exercise 1.2 we have that the twisted cubic has ideal $\gen{y-x^2,z-x^3}$. But the projective closure of the twisted cubic has ideal generated by $\gen{wy-x^2,wz-x^3,xy-zw}$. \qed
\lspace
\lspace

\textit{Note.} Lazy proof I know, but trust me when I say I did spent time thinking through this problem, and I think I get the important ideas. See the dm I sent to you for my main takeaways.

\subsection*{Question 2.10}

Skipped.

\subsection*{Question 2.11}

\noindent a)

Let $\set{p_i}$ denote a family of pairwise coprime linear polynomials.

Then $Y = \bigcap_i Z(p_i) \iff Y = Z(\bigcup_i p_i) \iff I(Y) = \gen{p_1, \dots, p_k}$. \qed

\noindent b)

Suppose that $Y \subseteq \mathbf{A}^n$ is a linear variety of dimension $r$. Then we have that $\text{ht}(I(Y)) = n-r$. But then if $I(Y)$ is generated by $m < n-r$ linear polynomials, by Krull's height theorem we have that $\text{ht}(I(Y)) \leq m$ which is a contradiction.

Now note that the zero set of a linear polynomial describes an affine subspace of $\mathbf{A}^n$. And as intersections of affine planes is still an affine plane, we get a correspondence between linear varieties and affine spaces.

Hence, by considering affine subvarieties of $Y$, we get that the dimension of $Y$, considered as an affine vector space over $k$, is less than or equal to the dimension of $Y$ as a topological space. Thus if we have a collection of linear polynomials that vanish on $Y$ that generated $I(Y)$, at most $n-r$ of them can be $k$-linearly independent. Hence I(Y) is minimally generated by $n-r$ linear polynomials. \qed

\noindent c)

By (b) above we have that $I(Y)$ and $I(Z)$ are generated by $n-r$ and $n-s$ polynomials respectively. Hence $I(Y \cap Z)$ is generated by $2n-r-s$ linear polynomials, such that $Y\cap Z$ has dimension at least $n - (2n-r-s) = r+s-n > 0$ by Question 1.9. Hence it cannot be empty. Furthermore by the reasoning given in (b), $Y\cap Z$ is a linear variety, and we are done. \qed
\lspace

\subsection*{Question 2.12}

\noindent a)

As the image is embedded within an integral domain, that the kernel is a prime ideal is clear.

Suppose that $f \in \ker \theta$, and let $m_i$ denote the sum of all homogeneous terms of degree $i$ in $f$, such that we get the decomposition $f = \sum_i m_i$. Then we get that $0 = \theta(f) = \sum_i \theta(m_i)$; and as the degree of two monomials under $\theta$ are the same iff they have the same degree we have that $\theta(m_i) =0$ for all $i$. Hence $\ker \theta$ is a prime homogeneous ideal, so $Z(\ker \theta)$ defines a projective variety in $\mathbf{P}^N$. \qed

\noindent b)

That $\rho_d(\mathbf{P}^n) \subseteq Z(\ker \theta)$ is clear: Suppose that $\vec p = (M_0(\vec a),\dots, M_N(\vec a)) \in \rho_d(\mathbf{P}^n)$ for some $\vec a \in \mathbf{P}^n$. Then for any $f \in \ker \theta$ such that $\theta(f) = f(M_0,\dots, M_N) = 0$, we automatically have that $f(\vec p) = f(M_0(\vec a),\dots, M_N(\vec a)) = 0$.


Now, notice that if $x_0^{n_0}\dots x_n^{n_n}$ denotes some arbitrary monomial with degree $d$, we have that $(x_0^{n_0}\dots x_n^{n_n})^d = (x_0^d)^{n_0} \cdots (x_n^d)^{n_n}$. Now let $\vec p \in Z(\ker \theta) \setminus \rho_d(\mathbf{P}^n)$. Consider the entries in $\vec p$ with index corresponding to the monomials $x_0^d, x_1^d, \dots, x_n^d$, and label them $p_0,\dots,p_n$. Then using the entries $p_i$ as a fixed reference, there must exist a coordinate $p'$ of $\vec p$, corresponding to the monomial $x_0^{n_0}\dots x_n^{n_n}$, such that $p'^d \neq (p_0^d)^{n_0} \cdots (p_n^d)^{n_n}$. But then $\vec p \not \in Z((x_0^{n_0}\dots x_n^{n_n})^d - (x_0^d)^{n_0} \cdots (x_n^d)^{n_n})$, which is a contradiction. Hence $\rho_d(\mathbf{P}^n) \supseteq Z(\ker \theta)$, and we get equality. \qed

\newpage
\noindent c)

That $\rho_d$ is bijective is clear. That it is a continuous map is also clear, as if $f \in k[y_0,\dots,y_N]$ then $\rho_d^{-1}(Z(f))$ is exactly $Z(\theta(f))$.

Now suppose $f \in k[x_0,\dots,x_n]$ is a homogeneous polynomial. We construct an ideal $\mathfrak a$ in $k[y_0,\dots,y_N]$ such that $\rho_d(Z(f)) = Z(\mathfrak a)$. This is sufficient, as every ideal is finitely generated, such that for an arbitrary closed set $Y$ we have that $I(Y) = \gen{f_1,\dots, f_m} = \gen{\bigcup_i f_i} = \bigcap_i Z(f_i)$.

Pick a $k \in \Z$ such that $d \mid k + \deg f$. Then for every $i$, $f_i := x_i^{k}f$ is a polynomial in the degree $d$ monomials and is hence contained in the image of $\theta$. So let $\mathfrak a := \theta^{-1}(\gen{f_i}_{i=0}^n)$. Then if $\rho_d(\vec p) \in Z(\mathfrak a)$, WLOG let $p_0$ be a nonzero coordinate of $\vec p$. Then $f_0(\vec p) = p_0^kf(\vec p) = 0$ and thus $f(\vec p) = 0$. Hence $\rho_d(Z(\mathfrak a)) = Z(f)$ as well. Thus $\rho_d$ is a closed map; and as it is a continuous bijection it is also a homeomorphism. \qed

\noindent d)

Label the 3-mononials in two variables as $M_0 := x^3, M_1:= x^2y, M_2:=xy^2, M_3 := y^3$. Then the 3-uple embedding of $\mathbf{P}^1$ in $\mathbf{P}^3$ is exactly the set $(M_0:M_1:M_2:M_3)$. As $\rho_d$ defines a homeomorphism this image is irreducible; and by (b) above image is closed. But then it contains the twisted cubic such that both have dimension $1$ so they are the same. \qed
\lspace

\subsection*{Question 2.13}

As $Z$ is a curve in $Y$, $\rho_2^{-1}(Z)$ is a curve in $\mathbf{P}^2$ via the homeomorphism constructed in 2.12. Hence write $\rho_2^{-1}(Z) = Z(f)$ for some homogeneous irreducible $f \in k[x,y,z]$. If $f$ has even degree then $f = \theta(g)$ for some $g \in k[x_0,\dots,x_5]$. Otherwise $f^2$ is an even degree homogeneous polynomial with $Z(f^2) = Z(f)$, such that there is a $g$ where $\theta(g) = f^2$. That $g$ is irreducible in the first case is clear. In the second case, note that any decomposition of $g$ corresponds to a decomposition of $f^2$; but the only one available is $f^2 = f \cdot f$, where $f$ cannot be in the image of $\theta$. Hence $g$ is always irreducible.

Thus, the bijection proven in 2.12 above gives us that $Z = Y \cap Z(g)$, where $Z(g)$ is the hypersurface as desired.
\lspace

\subsection*{Question 2.14}

Let the homogeneous coordinates of $\mathbf{P}^N$ be denoted by $(z_{ij})$, and consider the homomorphism $\theta: k[z_{ij}] \to k[x_0,\dots,x_r,y_0,\dots,y_s]$ given by $z_{ij} \mapsto x_iy_j$. Note that $\ker \theta$ is clearly a prime ideal as the image is contained in an integral domain. Next we show that $Z(\ker \theta) = \text{Im}(\psi)$. The proof proceeds as very similar to that in 2.12(b):


That $\text{Im}(\psi)) \subseteq Z(\ker \theta)$ is clear: Suppose that $\vec p = (p_{x_i}p_{y_j}) \in \text{Im}(\psi)$. Then for any $f \in \ker \theta$ such that $\theta(f) = f(x_iy_j) = 0$, we automatically have that $f(\vec p) = f(p_{x_i}p_{y_j}) = 0$.

Suppose that $\vec p \in Z(\ker \theta) \setminus \text{Im}(\psi)$. Then, using the entries $p_{0j}$ and $p_{i0}$ of $\vec p$ as fixed reference, there exists some $p_{kl}$ such that $p_{00}p_{kl} \neq p_{0l}p_{k0}$. But then $\vec p \not \in Z(z_{00}z_{kl} - z_{0l}z_{k0})$, where $\theta(z_{00}z_{kl} - z_{0l}z_{k0}) = x_0y_0 x_ky_l - x_oy_lx_ky_0 = 0$; which is a contradiction. Hence $\text{Im}(\psi) \supseteq Z(\ker \theta)$, and we get equality. \qed
\lspace

\subsection*{Question 2.15}

\noindent a)

Define the embedding $\phi: \mathbf{P}^1 \times \mathbf{P}^1 \hookrightarrow \mathbf{P}^3$ where $(x_0:x_1),(y_0;y_1) \mapsto (x_0y_0:x_0y_1:x_1y_0:x_1y_1)$. Then $x_0y_1x_1y_0 - x_0y_0x_1y_1 = 0$, and a similar proof as done in 2.14 yields that $\text{Im}(\phi) = Z(xy-zw)$. \qed

\noindent b)

In both copies of $\mathbf{P}^1$, notice that we can WLOG assume the first entry is nonzero such that every point in the space has the form $(1:t)$ where $t$ ranges over every value in $k$.

If $t$ is fixed for the first component, then its image under the Segre embedding corresponds to the set $(y_0:y_1:ty_0:ty_1)$ where $(y_0:y_1)$ ranges over $\mathbf{P}^1$. This is a linear projective variety, with ideal given by $\gen{y-tw,z-tx}$. It can be seen that the ideal has height $2$, so the projective variety has dimension $1$ as desired. The construction for the second component is similar, and these give the families $\set{L_t}$ and $\set{M_t}$.

Now if $u \neq t$, notice that the projective variety $L_t \cap L_u$ has ideal given by $\gen{y-tw, z-tx,y-uw, z-ux}$. But then these are $4$ algebraically independent linear polynomials in $k[w,x,y,z]$ such that they generate the entire ring. Hence by the homogeneous Nullstellensatz the intersection is empty. The proof for $M_t \cap M_y = \emptyset$ is analogous.

Now consider $L_u \cap M_t$. The ideal of the intersection is given by $\gen{y-tw, z-tx,w-ux, z-uy}$. If both $t,u=0$ then it is clear that the coordinate ring is isomorphic to $k$. Otherwise suppose that $t \neq 0$, such that in the coordinate ring we have that $w \equiv ux \equiv \frac{u}{t}z \equiv \frac{u^2}{t}y \equiv u^2 w$, such that $w \equiv 0$. But then this proof applies for $x,y,z$ as well such that the coordinate ring is also isomorphic to $k$. Hence the ideal is a maximal ideal, and by the homogeneous Nullstellensatz is exactly given by the kernel of the evaluation homomorphism at some point $\vec p \in \mathbf{P}^3$. In other words, $L_u \cap M_t = Z(\gen{y-tw, z-tx,w-ux, z-uy}) = \set{\vec p}$. \qed
\lspace

\subsection*{Question 2.16}

\noindent a)

Let $Q_1 := Z(x^2-yw)$ and $Q_2 := (xy-zw)$, both considered as projective varieties in $\mathbf{P}^3$.

Then the points in the intersection $Q_1 \cap Q_2$ are given by $(w:x:\frac{x^2}{w}:\frac{x^3}{w^2}) = (w^3:w^2x:wx^2:x^3)$ if $w \neq 0$, which corresponds to a twisted cubic, and $(0:0:y:z)$ if $w=0$, which is a line. \qed

\noindent b)

Let $C := Z(x^2-yz)$ and $L:=Z(y)$, both considered as projective varieties in $\mathbf{P}^3$.

Then the points in the intersection are given by $(0:0:z)$, i.e. is a single point. But then the sum of the ideals is $\gen{x^2,y}$, which is not a prime ideal and hence not the minimal ideal corresponding to $C \cap L$. \qed
\lspace

\subsection*{Question 2.17}

\noindent a)

Suppose that $Y \subseteq \mathbf{P}^n$ is a projective ideal such that $\mathfrak a:= I(Y)$ can be generated by $q$ elements. Then Krull's height theorem gives us that $\text{ht}(\mathfrak a) \leq q$, such that by the result of 2.6 it immediately follows that $\dim(Y) \geq n-q$. \qed

\noindent b)

Suppose that $Y$ is a strict complete intersection such that $I(Y) = \gen{f_1,\dots,f_q}$. Then $Y \supseteq \bigcap_i Z(f_i) \supseteq Z(\gen{f_1,\dots,f_q}) = Y$, such that $Y$ is also a set-theoretic complete intersection. \qed

\noindent c)

Skipped.

\noindent d)

Nope.

\end{document}